% © 2026 Ing. David Jaroš — CC BY-NC-ND 4.0
%
% This work is licensed under a Creative Commons Attribution-NonCommercial-NoDerivatives
% 4.0 International License (CC BY-NC-ND 4.0).
%
% License History: Earlier drafts (up to v0.3) were released under CC BY 4.0.
% From v0.4 onward, all material is released under CC BY-NC-ND 4.0 to protect
% the integrity of the theoretical work during ongoing academic development.
%
% See LICENSE.md for full license text.

\documentclass[12pt,a4paper]{article}
\usepackage[a4paper,margin=2.5cm]{geometry}
\usepackage{amsmath,amssymb,amsthm,booktabs,tcolorbox}
\tcbuselibrary{skins,breakable}
\newtcolorbox{statusbox}[1][]{colback=gray!8!white,colframe=black!60,arc=3pt,boxrule=1pt,title=Status Summary,#1}

\title{Gap G137-B: Chowla--Selberg coefficient $B$ on $\mathbb{Z}^2$}
\author{Ing.\ David Jaroš}
\date{2026-05-12}

\begin{document}
\maketitle

\section{Goal}
Derive the coefficient bridge in
\[
V_{\mathrm{eff}}(n)=n^2-B\,n\ln n
\]
from first principles using the zeta-regularised determinant chain for the square torus and its mass deformation.

\section{Step 1: Chowla--Selberg / Epstein baseline [STD]}
For the square lattice,
\[
Z_{\mathbb{Z}^2}(s)=4\,\zeta(s)L(s,\chi_{-4}).
\]
Hence
\[
Z'_{\mathbb{Z}^2}(0)=4\big(\zeta'(0)L(0,\chi_{-4})+\zeta(0)L'(0,\chi_{-4})\big).
\]
Using
\(\zeta(0)=-\tfrac12\),
\(\zeta'(0)=-\tfrac12\ln(2\pi)\),
\(L(0,\chi_{-4})=\tfrac12\)
and the functional equation for
\(L(s,\chi_{-4})\), one gets an explicit form involving
\(\psi(\tfrac14),\psi(\tfrac34)\), thus a logarithmic dependence on
\(\Gamma(\tfrac14)\).

\section{Step 2: $\Gamma(1/4)$, $\eta(i)$ and $\vartheta_3(0|i)$ [STD/L0]}
Classical CM identities:
\[
\eta(i)=\frac{\Gamma(1/4)}{2\pi^{3/4}},
\qquad
\vartheta_3(0|i)=\frac{\pi^{1/4}}{\Gamma(3/4)},
\qquad
\frac{\vartheta_3(0|i)}{\eta(i)}=\sqrt2.
\]
Therefore any logarithmic insertion built from
\(Z'_{\mathbb{Z}^2}(0)\)
can be re-expressed algebraically in terms of
\(\ln\eta(i)\)
and
\(\ln\vartheta_3(0|i)\).

\section{Step 3: mass derivative channel for $n\ln n$}
For massive deformation of the torus determinant,
\[
\frac{\partial}{\partial m^2}Z'_{T^2}(0;m^2)\Big|_{m\to0}
= -Z_{T^2}(1)_{\mathrm{ren}},
\]
and divisor-sum asymptotics of the $T^2$ winding channel provide the
\(n\ln n\) structure [L1 asymptotically].
The unresolved step is the exact insertion map from this derivative to the closed coefficient
\(B\)
in the effective potential normalisation used by UBT.

\section{Step 4: target expression and current closure status}
Target expression:
\[
B_{\mathrm{target}}=12^{3/2}\,2^{1/8}\,\vartheta_3(0|i)^{1/4}.
\]
Numerically this remains an observation-level near-match unless the determinant-to-coefficient map is closed analytically inside the same renormalisation scheme.

\begin{statusbox}
Chowla-Selberg $Z'_{\mathbb{Z}^2}(0)$ contains $\ln\Gamma(1/4)$: \textbf{[L1]}\\
$\Gamma(1/4)\to\eta(i)\to\vartheta_3(0|i)$ algebraically: \textbf{[L0/STD]}\\
$B=f(\vartheta_3(0|i))$ analytically in the UBT normalisation: \textbf{[OPEN --- exact determinant-to-coefficient insertion not closed]}\\
Gap G137-B: \textbf{[NARROWED --- analytic modular chain fixed; coefficient insertion step remains]}\\
Alpha je odvozena: \textbf{NE (unconditional dokud [L1] není dosaženo)}
\end{statusbox}

\end{document}
